\renewcommand*{\authorOfNextLines}{Helmut Quinz}
\begin{quotation}
	\textit{Das Zitat (lateinisch citatum \gqq{Angeführtes, Aufgerufenes} zu lat. citāre \gqq{in Bewegung setzen, vorladen}, vgl. \gqq{jemanden vor Gericht zitieren}) 
	ist eine wörtlich oder inhaltlich übernommene Stelle aus einem Text oder ein Hinweis auf eine bestimmte Textstelle. 
	Auch Inhalte aus anderen Medien können übernommen werden: Es gibt Bild-, Musik-, Film- und Architekturzitate.}
	(\cite{dewiki:236507884})
\end{quotation}
Ein wesentliches Prinzip wissenschaftlichen Arbeitens ist die Nachvollziehbarkeit der in einem Dokuments getätigten Aussagen. 
Werden in einer (vor)wissenschaftlichen Arbeit fremde Quellen den eigenen Aussagen zugrunde gelegt (also \gqq{zitiert}), so sind diese Quellen vollständig und korrekt anzugeben!(vgl. \cite{Bergler2020})\\
Quellen solcher Zitate können sein:
\begin{itemize}
	\item Texte (Bücher, Fachzeitschriften, Produktinformationen, Firmenunterlagen etc.)
	\item Filme, Videosequenzen
	\item Radiosendungen→ Unterrichtsinhalte
	\item Grafiken (Diagramme, Tabellen etc.)
	\item Informationen aus dem Internet
	\item persönliche Mitteilungen, z.B. externer Fachexperten
\end{itemize}
\par

Unterschieden wird zwischen dem
\begin{itemize}
	\item direktem Zitat in den Formen eines
	\begin{itemize}
		\item Kurzzitats oder eines 
		\item Blockzitats
	\end{itemize}
		und dem 
	\item indirekten (sinngemäßen) Zitat
\end{itemize}
\FloatBarrier